\chapter{Skew and Term-Structure Relative Value}\label{ch:s2:skew-and-term-structure-relative-value}

From September 2009 to September 2026 Cboe's three-month volatility index stood on average 1.93 points above its one-month VIX. On 7.6\% of days the
order was reversed, and on 12 March 2020 the one-month index stood 18.23 points above the three-month. Two options on the same index, one a month out
and one three months out, can disagree about the same future. Relative-value volatility traders bet on the disagreement, or on the shape of the smile,
rather than on the level. This chapter builds a synthetic surface whose skew and slope carry planted mispricings that fade, and trades them with
vega-neutral, delta-hedged options. The mispricings are real: the parts of the P\&L that come from them have Sharpe ratios of 1.0 and 1.5. The trades
earn 0.75 and 0.64, and most of what they earn comes from something else. The build is \texttt{firm.volrv}.

\section{Trading the smile}

An index's implied volatility is higher for low strikes than for high ones: the skew (Book~1, chapter~25). Cboe publishes a measure of it. The SKEW
index is $100 - 10S$, where $S$ is the price of S\&P 500 skewness computed from a portfolio of options; it rises as the risk-neutral distribution's
left tail grows. From January 1990 to September 2026 it averaged 123.2, with 5\% and 95\% quantiles of 110.0 and 147.3. Its low was 101.23 on 19 March
1991 and its high 183.12 on 18 February 2025. Cboe's white paper noted the low correlation between variations in SKEW and VIX. On the history since
1990, daily changes in the two correlate at $-0.13$, and their levels at $-0.17$.

Bakshi, Kapadia and Madan compared the index with its members, using options on the S\&P 100 and 30 of its stocks. Individual stocks' risk-neutral
distributions were far less negatively skewed than the index's. They traced skewness to risk aversion and decomposed each stock's into a systematic
and an idiosyncratic part.

\begin{definition}[Skew trade]\label{def:s2:skew-and-term-structure-relative-value:skew}
A \emph{skew trade}\index{skew trade} buys options at one strike and sells options at another on the same underlying and expiry, sized so that the
position has no net vega and is delta-hedged, in order to profit from a change in the difference between their implied volatilities rather than from
the level.
\end{definition}

The usual vehicle is the risk reversal (Book~2, chapter~19): long an out-of-the-money call, short an out-of-the-money put, or the reverse. A trader
who thinks the put wing is too dear sells the put and buys the call. The trade has a second exposure hidden in it. A surface moves with spot as well as
with time: under a sticky-strike rule (Book~5, chapter~7) each strike keeps its vol as spot moves, and under a sticky-delta rule the smile moves with
spot. When the index falls, the short put comes nearer the money and its vega grows. The position that had no vega now has some, and it is short.

\subsection*{The synthetic surface}

\texttt{firm.volrv} puts a daily surface on the synthetic index of chapter~1. The fair one-month and three-month at-the-money vols come from the
variance dynamics, premium included. The fair skew coefficient, the slope of the smile per unit of standardised moneyness $x = \ln(K/S)/(\sigma\sqrt\tau)$,
is $-0.12$ at the average level. It steepens as the level rises: 0.4 of a unit of skew for each unit of one-month vol. On top of the fair surface the
market adds two planted noises: one on the skew coefficient (standard deviation 0.06) and one on the three-month vol (1.5 vol points). Each is an
autoregression with a daily persistence of 0.98, a half-life of 34 days. They stand for demand that pushes one part of the surface and then fades. The
smile is $\sigma(x) = \sigma_{\mathrm{atm}}(1 + \text{skew}\cdot x + 0.02x^2)$.

A trader does not see the noises. \texttt{features} extracts what can be seen: the level (the one-month vol), the slope (three-month minus one-month) and
the skew. It then regresses slope and skew on the level, using only data up to each day, and keeps the residuals. The raw skew correlates with the
level at $-0.63$ and the raw slope at $-0.83$. The full-sample fits give $-0.46$ for the skew's level beta (the planted value is $-0.4$) and $-0.20$ for
the slope's. The residuals recover the planted noises well: they correlate with them at 0.97 for the skew and 0.92 for the slope.

\section{Trading the term structure}

\begin{definition}[Term-structure trade]\label{def:s2:skew-and-term-structure-relative-value:term}
A \emph{term-structure trade}\index{term-structure trade} buys volatility at one maturity and sells it at another on the same underlying, sized so
that the position has no net vega and is delta-hedged, in order to profit from a change in the slope of the term structure rather than from the level.
\end{definition}

The calendar spread is the standard form: long three-month, short one-month at-the-money options with equal vegas. Its shape follows from how vol
mean-reverts. When vol is high, the market expects it to fall, so near-dated vol stands above far-dated vol and the term structure inverts. When vol is
low, the order reverses. A slope is therefore, in large part, a level in disguise. On the real indices (\cref{fig:s2:skew-and-term-structure-relative-value:slope})
the slope VIX3M minus VIX correlates with VIX at $-0.54$. It averaged $-7.98$ points through March 2020. It is also quick to revert: its daily
persistence is 0.90, a half-life of 6.5 trading days.

\begin{omfigure}
\begin{tikzpicture}
\begin{axis}[width=11cm,height=5.4cm,xlabel={\small year},ylabel={\small VIX3M $-$ VIX (points)},tick label style={font=\footnotesize},
  xmin=2009.5,xmax=2027,ymin=-9,ymax=7,grid=major,grid style={gray!25},table/col sep=comma,/pgf/number format/1000 sep={}]
\addplot[omThm,thick] table[x=year,y=slope]{figdata/strategies-2/03-skew-and-term-structure-relative-value/slope.csv};
\addplot[black,thin,domain=2009.5:2027] {0};
\end{axis}
\end{tikzpicture}

\omcaption{The slope of the VIX term structure: Cboe's three-month volatility index (VIX3M) minus VIX, monthly means, September 2009 to September 2026.
Below zero the curve is inverted: near-dated vol stands above far-dated vol. Derived from Cboe's index histories; the raw series are not redistributed.
Data: \texttt{s2\_fetch\_term}.}\label{fig:s2:skew-and-term-structure-relative-value:slope}
\end{omfigure}

Johnson studied what the slope says. The shape of the VIX term structure carries information about the price of variance risk, not about expected
changes in VIX, which rejects the expectations hypothesis. Its second principal component, a slope, predicts the excess returns of synthetic S\&P 500
variance swaps, VIX futures and S\&P 500 straddles at all maturities. A trader can read this two ways. It is a timing signal for the level trades of
chapter~1: sell variance when the slope says the premium is high. It is also a warning for a calendar spread, which trades the slope itself.

\section{Across related underlyings}

\begin{definition}[Cross-underlying volatility spread]\label{def:s2:skew-and-term-structure-relative-value:cross}
A \emph{cross-underlying volatility spread}\index{cross-underlying volatility spread} buys implied volatility on one underlying and sells it on a
related one, such as two indices, an index and an exchange-traded fund on it, or an index and its members, betting that the ratio or difference of the
two vols returns to its usual range.
\end{definition}

Cboe publishes volatility indices for the Russell 2000 (RVX) and the Nasdaq-100 (VXN). Since September 2009, RVX has stood on average 5.2 points above
VIX, with a standard deviation of 2.3 points and a half-life of 15.7 trading days. VXN has stood 3.0 points above, with a standard deviation of 2.3 and
a half-life of 15.3 days. The spreads revert more slowly than the VIX slope, and nothing in them is the level of one market in disguise. That makes them
the cleanest relative-value signal of this chapter. It also brings the least familiar risk. The two underlyings can part company for a reason (a
sector's crash, a change in an index's membership), and the spread then does not come back.

Nobody trades the volatility indices themselves: the trade is in options on each underlying. It carries each underlying's delta hedge, each option
market's spread, and the chance that one surface is quoted less often than the other. Bakshi, Kapadia and Madan's result is the basis of the most
common cross-underlying trade, index skew against single-stock skew. Its risk is the correlation of chapter~2.

\section{The books and their attribution}

Each month the book strikes a unit trade scaled by minus the signal, an expanding z-score clipped at 2: it sells what is rich. The trade is a risk
reversal one standardised unit either side of the money, or a one-month against three-month calendar. Each leg is sized so that a one-percent rise in
its vol is worth one unit (0.18 vol points at a vol of 18). Every leg is delta-hedged daily and held to the month's end. Opening a leg costs one
percent of its vol, the chapter's own assumption rather than a sourced spread. The P\&L is attributed each day by repricing in order
(\cref{lst:s2:skew-and-term-structure-relative-value:pnl}): spot and time (delta, gamma, theta, on the day's surface), then the level (the fair surface
moving to the next day's), then the skew noise, then the slope noise. The order matters: the level is moved first, and the skew and slope buckets hold
only the planted noises. Four books run over 227 months after a year's warm-up: each trade, signalled by the raw feature or by its residual on the
level.

\begin{center}
\small
\begin{tabular}{@{}lccccccc@{}}
\toprule
& & & & \multicolumn{4}{c}{share of P\&L (\%)}\\
\cmidrule(l){5-8}
227 months & Sharpe & skewness & worst (sd) & spot, time & level & skew or slope & cost\\
\midrule
skew, raw & 1.13 & 1.91 & $-4.6$ & 94 & 5 & 9 & $-8$\\
skew, residual & 0.75 & 1.48 & $-4.5$ & 73 & 24 & 16 & $-13$\\
calendar, raw & 0.07 & 4.74 & $-1.7$ & $-31$ & 160 & 155 & $-186$\\
calendar, residual & 0.64 & 1.43 & $-2.4$ & 97 & $-16$ & 51 & $-32$\\
\bottomrule
\end{tabular}
\end{center}

The planted effects are there. Taken alone, the skew bucket of the residual skew book has a Sharpe ratio of 1.01 and the slope bucket of the residual
calendar 1.51. The raw books catch less of them: 0.95 for the raw skew book's skew bucket and 0.84 for the raw calendar's slope bucket. But these
buckets are small next to the rest. The planted skew noise explains 16\% of the residual skew book's P\&L. The skew book signalled on the raw feature
has the \emph{highest} Sharpe ratio of the four, and only 9\% of its P\&L comes from the skew moving back. The rest is spot and time. The raw skew is
steep when vol is high, so this book sells the put wing when vol is high. What it collects is chapter~1's premium, with a risk reversal's paperwork.

The raw calendar is the other case. Its slope signal is 83\% level, so its level bucket (160\% of the P\&L) and its slope bucket (155\%) are both
large, and costs take the difference. The residual calendar keeps the slope bucket and sheds most of the level. Half its P\&L is the planted effect.
The other half is spot and time, and part of that is the mispricing too: an option bought cheap earns more from gamma than it pays in theta while it is
held. Across months the spot-and-time bucket correlates with the size of the mispricing held at 0.25 for the skew book and 0.14 for the calendar. The
rest is gamma noise, and it dominates the month-to-month variance (\cref{fig:s2:skew-and-term-structure-relative-value:calendar}).

\begin{omfigure}
\begin{tikzpicture}
\begin{axis}[width=11cm,height=5.4cm,xlabel={\small year},ylabel={\small cum.\ P\&L (sd units)},tick label style={font=\footnotesize},
  xmin=1,xmax=20,ymin=-15,ymax=50,grid=major,grid style={gray!25},table/col sep=comma,
  legend style={font=\scriptsize,at={(0.5,-0.32)},anchor=north},legend columns=4]
\addplot[black,thick] table[x=year,y=total]{figdata/strategies-2/03-skew-and-term-structure-relative-value/calendar.csv};
\addplot[omDef,thick] table[x=year,y=spot_time]{figdata/strategies-2/03-skew-and-term-structure-relative-value/calendar.csv};
\addplot[omThm,thick] table[x=year,y=slope]{figdata/strategies-2/03-skew-and-term-structure-relative-value/calendar.csv};
\addplot[gray,thick,dashed] table[x=year,y=level]{figdata/strategies-2/03-skew-and-term-structure-relative-value/calendar.csv};
\legend{total,spot and time,slope noise,level}
\end{axis}
\end{tikzpicture}

\omcaption{The residual calendar book's P\&L and three of its buckets, cumulated in units of the book's monthly standard deviation (costs, not shown,
make up the difference). The planted slope noise pays steadily. Spot and time is larger and far noisier. Data:
\texttt{s2\_volrv.books}.}\label{fig:s2:skew-and-term-structure-relative-value:calendar}
\end{omfigure}

One seed can flatter a book. Over eight seeds of the market and the surface, the mean Sharpe ratios are 0.90 for the raw skew book (range 0.47 to
1.26), 0.39 for the residual skew book (0.12 to 0.75), $-0.07$ for the raw calendar ($-0.43$ to 0.29) and 0.38 for the residual calendar ($-0.10$ to
0.75). The chapter's seed is on the lucky side for both residual books.

\section{Risk: what a relative-value vol book is really short}

A book with no net vega and no net delta at inception is not a book without risk. Held every month without a signal, one unit of each trade shows
what it carries. Long the call wing and short the put earns 12.0 units a month from spot and time and loses 15.9 to the level: $-6.2$ a month in
all. Its short put gains vega as the index falls, when vol rises. Long the calendar earns 9.6 units a month from spot and time, since the short
one-month leg collects the variance premium through theta, and loses 3.2 to the level: 4.7 a month in all.

The worst month of both skew books began on day 1\,470 (year 5.8). Both books were at their limit: two units long the call wing and short the put. The
crash came nine days after the month ended. Inside the month the index rose 4.5\% while the one-month vol fell from 15.2\% to 9.7\%. The long call
rose to its strike and its vega grew just as the vol it was long collapsed. The month cost 289 units: 119 from spot and time (the call's theta in a
calm rally) and 157 from the level, 4.5 of the book's monthly standard deviations. The residual signal had no way to see it coming: the skew noise was
steep that month and said to sell the put wing. A relative-value vol book is short whatever its legs do not match. For a risk reversal that is a joint
move of spot and vol in either direction: a crash when it is short the put wing, a calm rally with falling vol when it is long the call. For a calendar
it is gamma in the near month, which the three-month leg does not cover.

Three rules follow. Attribute before believing a Sharpe ratio: a \omterm{def:s2:skew-and-term-structure-relative-value:skew}{skew trade} whose skew bucket is a tenth of its P\&L is a different trade. Signal on
residuals to the level, or the book is a level trade with extra costs. Stress each book on joint moves of spot, vol and skew in both directions, not on
one of them at a time.

\section{Strategy files}

\begin{strategyfile}{Skew flattener}\label{strat:s2:skew-and-term-structure-relative-value:skew}
\sfield{Who pays you, and why} Hedgers who bid index puts beyond what the surface's other points justify; the payment is the part of the put wing's
richness that fades.
\sfield{Instruments and venues} Listed index options, out-of-the-money puts and calls of one expiry.
\sfield{Signal} The skew's residual after its dependence on the level, against its history.
\sfield{Sizing and execution} Vega-neutral risk reversals, delta-hedged daily; sized by the crash-month loss, not by the Sharpe ratio.
\sfield{Costs} Two option spreads a roll and the hedging; the put wing is the less liquid.
\sfield{How it dies} A crash while short the put wing; a skew that steepens with the level and never comes back.
\sfield{Horizon, capacity, infrastructure} Weeks to months; a daily surface fit and a hedging engine.
\sfield{Backtest honestly} Attribute the P\&L; test the raw and residual signals; include crash months.
\sfield{Sources} Bakshi, Kapadia and Madan (2003) on index skew; this chapter: 0.75 on one seed, 0.39 on average over eight, with 16\% of the P\&L
from the skew.
\end{strategyfile}

\begin{strategyfile}{Calendar volatility spread}\label{strat:s2:skew-and-term-structure-relative-value:calendar}
\sfield{Who pays you, and why} Demand concentrated at one maturity (hedgers in the near month, structured-product issuers further out) that fades.
\sfield{Instruments and venues} Listed index options at two expiries; variance swaps or VIX futures at two tenors.
\sfield{Signal} The slope's residual after its dependence on the level.
\sfield{Sizing and execution} Equal vegas; delta-hedged; near leg rolled monthly.
\sfield{Costs} Two spreads a roll; the near leg's hedging.
\sfield{How it dies} A vol spike, when the short near leg's gamma loses; a slope signal that is really the level.
\sfield{Horizon, capacity, infrastructure} Weeks to months; a term-structure fit.
\sfield{Backtest honestly} Signal on residuals; attribute to level and slope; include inversions.
\sfield{Sources} Johnson (2017) on the VIX slope and variance risk premia; this chapter: 0.64 on one seed, 0.38 on average.
\end{strategyfile}

\begin{strategyfile}{Index-versus-ETF volatility spread}\label{strat:s2:skew-and-term-structure-relative-value:etf}
\sfield{Who pays you, and why} Flows that hit one of two option markets on nearly the same underlying, such as index options against options on an
exchange-traded fund that tracks the index.
\sfield{Instruments and venues} Index options against the fund's options; the fund's tracking and dividends set the fair spread.
\sfield{Signal} The implied-vol difference against the tracking-adjusted fair difference.
\sfield{Sizing and execution} Vega-neutral, delta-hedged in each underlying.
\sfield{Costs} Two option markets' spreads; different settlement and exercise styles.
\sfield{How it dies} Differences that are fair (exercise style, settlement, dividends, taxes) mistaken for mispricing.
\sfield{Horizon, capacity, infrastructure} Days to weeks; both surfaces live.
\sfield{Backtest honestly} Model settlement and exercise differences before calling a gap a signal.
\sfield{Sources} No performance figure verified.
\end{strategyfile}

\begin{strategyfile}{Cross-market volatility spread}\label{strat:s2:skew-and-term-structure-relative-value:cross}
\sfield{Who pays you, and why} Demand for protection on one market that pushes its vol away from a related market's.
\sfield{Instruments and venues} Options or volatility futures on two indices (small caps against large caps, one region against another).
\sfield{Signal} The vol spread or ratio against its history; Cboe's RVX and VXN against VIX revert with half-lives of about 15 days.
\sfield{Sizing and execution} Vega-neutral or beta-weighted legs.
\sfield{Costs} Two markets' spreads and hedges.
\sfield{How it dies} A regime change in one market: the spread moves and stays.
\sfield{Horizon, capacity, infrastructure} Weeks; two surfaces.
\sfield{Backtest honestly} Test whether the spread reverts after a shock or trends.
\sfield{Sources} Cboe's RVX, VXN and VIX histories (statistics in this chapter); no performance figure verified.
\end{strategyfile}

\section{Tutorial: same future, two prices}\label{tut:s2:skew-and-term-structure-relative-value:tut}

\begin{tutorial}
\textbf{Goal.} Build the synthetic surface, extract its features, run the four books and attribute their P\&L; read the real term structure, SKEW and
cross-index spreads. \textbf{End state:} the table and the two figures.
\begin{tutsteps}
\item \textbf{The trade and its attribution}.
\omcode{code/firm/volrv/firm_volrv.py}{104}{137}{A unit trade struck monthly, delta-hedged daily, attributed by sequential repricing.}\label{lst:s2:skew-and-term-structure-relative-value:pnl}
\item \textbf{Signals and books}.
\omcode{code/strategies-2/03-skew-and-term-structure-relative-value/python/s2_volrv.py}{35}{64}{Expanding z-scores and the four books.}\label{lst:s2:skew-and-term-structure-relative-value:books}
\item \textbf{Run} \texttt{s2\_fetch\_term.py} once, then \texttt{table()}, \texttt{bucket\_sr()}, \texttt{seeds()}, \texttt{static()} and
\texttt{fig\_volrv.py}.
\end{tutsteps}
\textbf{What to change next.} Hedge the risk reversal's crash exposure with a small long put further out and measure what it costs; switch the
surface's rule to sticky delta and compare the attributions; add a third maturity and trade the curvature of the term structure.
\end{tutorial}

\section{Build: volatility relative value}\label{bld:s2:skew-and-term-structure-relative-value:volrv}

\begin{build}
\bfield{Purpose} A daily surface with planted skew and slope mispricings; surface features and their residuals on the level; vega-neutral risk
reversals and calendars with P\&L attribution.
\bfield{Interface} \texttt{SurfaceConfig(\dots)}, \texttt{surface\_path(r, v, cfg, scfg)}, \texttt{features(s, warmup)}, \texttt{trade\_pnl(r, s,
kind, pos, cfg, scfg)}, \texttt{attribution(b)}.
\bfield{Rules} Residuals from expanding regressions only; legs sized by relative vega at inception; attribution in a fixed order (spot and time,
level, skew, slope) so that the buckets add up to the total.
\bfield{Acceptance tests} \texttt{code/firm/volrv/tests/}: a linear level effect removed exactly; a still surface gives no level, skew or slope
P\&L and buckets that add up; a risk reversal gains when the skew flattens and loses when reversed.
\bfield{Stretch} Sticky-delta dynamics; curvature trades; two underlyings.
\end{build}

\begin{omsources}
\item G.~Bakshi, N.~Kapadia and D.~Madan, ``Stock return characteristics, skew laws, and the differential pricing of individual equity options'',
\emph{Review of Financial Studies} 16(1), 2003.
\item T.~L.~Johnson, ``Risk premia and the VIX term structure'', \emph{Journal of Financial and Quantitative Analysis} 52(6), 2017.
\item Cboe, SKEW Index white paper, 2011.
\item Cboe Global Markets, daily histories of VIX, VIX3M, SKEW, RVX and VXN.
\end{omsources}

\section{Exercises}

\begin{exercise}[$\star$]\label{exo:s2:skew-and-term-structure-relative-value:1}
The price of S\&P 500 skewness is $S = -2.3$. What is the SKEW index?
\end{exercise}

\begin{exercise}[$\star$]\label{exo:s2:skew-and-term-structure-relative-value:2}
A leg is sized so that a one-percent rise in its vol is worth one unit. Its vol is 20\% and falls to 19.6\%. What does a short position in the leg make?
\end{exercise}

\begin{exercise}[$\star$]\label{exo:s2:skew-and-term-structure-relative-value:3}
An autoregression has daily persistence 0.98. What is its half-life in trading days?
\end{exercise}

\begin{exercise}[$\star\star$]\label{exo:s2:skew-and-term-structure-relative-value:4}
Why is a calendar spread signalled on the raw slope largely a trade on the level of vol?
\end{exercise}

\begin{exercise}[$\star\star$]\label{exo:s2:skew-and-term-structure-relative-value:5}
Johnson found that the VIX slope predicts variance-swap, VIX-future and straddle returns. How would you use that finding in a level trade, and what does
it warn a calendar trader about?
\end{exercise}

\begin{exercise}[$\star\star$]\label{exo:s2:skew-and-term-structure-relative-value:6}
Individual stocks' risk-neutral distributions are less negatively skewed than the index's. Design a trade on that difference and name its main risk.
\end{exercise}

\begin{exercise}[$\star\star\star$]\label{exo:s2:skew-and-term-structure-relative-value:7}
\emph{Coding.} Run \texttt{table(0, 0.0)}, which switches off both planted noises. What happens to each book, and what does it tell you about the raw
skew book?
\end{exercise}

\begin{exercise}[$\star\star\star$]\label{exo:s2:skew-and-term-structure-relative-value:8}
\emph{Find the flaw.} ``Our skew book has a Sharpe ratio of 1.13 and the skew mean-reverts, so we are being paid for skew mean reversion.''
\end{exercise}

\section{Problem: Same Future, Two Prices}

\begin{problem}[{Weekend problem --- relative value on a surface}]\label{pb:s2:skew-and-term-structure-relative-value:1}
The chapter's synthetic surface, Cboe's indices and the public record.

\medskip\noindent\textbf{Part I --- The surface.}
\begin{enumerate}
\item Define a \omterm{def:s2:skew-and-term-structure-relative-value:skew}{skew trade} and a \omterm{def:s2:skew-and-term-structure-relative-value:term}{term-structure trade}.
\item How is Cboe's SKEW index defined, and what were its statistics?
\item What did Bakshi, Kapadia and Madan find about index and stock skews?
\item How does the synthetic surface's fair skew depend on the level, and what are its planted noises?
\end{enumerate}

\medskip\noindent\textbf{Part II --- The real term structure.}
\begin{enumerate}[resume]
\item Give the VIX slope's average, how often it was inverted, and its extreme.
\item Why does the slope correlate with the level?
\item What did Johnson find?
\item Give the RVX and VXN spreads to VIX and their half-lives.
\end{enumerate}

\medskip\noindent\textbf{Part III --- The books.}
\begin{enumerate}[resume]
\item How are the features and their residuals extracted?
\item How is each leg sized, and in what units is the P\&L?
\item How is the P\&L attributed, and why does the order matter?
\item Give the four books' Sharpe ratios and their attributions.
\end{enumerate}

\medskip\noindent\textbf{Part IV --- The verdict.}
\begin{enumerate}[resume]
\item State the \emph{named result}: the skew and calendar trades' Sharpe ratios and the share of their P\&L that is really something else.
\item What do the planted effects earn alone?
\item What did the eight seeds show?
\item What does each trade carry when held without a signal?
\item What happened in the worst month of the skew books?
\item What happens when the planted noises are switched off?
\item How would you backtest a skew book honestly?
\item In one sentence: what is a relative-value vol book short?
\end{enumerate}
\end{problem}

\section{Interview questions}

\begin{interviewq}[$\star$ \iqroles{trader}]\label{iq:s2:skew-and-term-structure-relative-value:1}
What is a risk reversal, and what does its price tell you?
\end{interviewq}

\begin{interviewq}[$\star\star$ \iqroles{trader}]\label{iq:s2:skew-and-term-structure-relative-value:2}
You are long a vega-neutral calendar. The market sells off sharply. What happens to your P\&L?
\end{interviewq}

\begin{interviewq}[$\star\star$ \iqroles{researcher}]\label{iq:s2:skew-and-term-structure-relative-value:3}
How would you tell whether a skew signal is a level signal in disguise?
\end{interviewq}

\begin{interviewq}[$\star\star$ \iqroles{risk}]\label{iq:s2:skew-and-term-structure-relative-value:4}
A relative-value vol book reports zero vega and zero delta. What risks remain?
\end{interviewq}

\begin{interviewq}[$\star\star$ \iqroles{developer}]\label{iq:s2:skew-and-term-structure-relative-value:5}
Design a daily P\&L attribution for an options book. What has to be true for the buckets to add up?
\end{interviewq}

\begin{interviewq}[$\star\star\star$ \iqroles{researcher}]\label{iq:s2:skew-and-term-structure-relative-value:6}
Under Black--Scholes, show that a delta-hedged option's P\&L over a short interval is approximately $\tfrac12\Gamma S^2(\sigma_r^2 -
\sigma_i^2)\,\mathrm{d}t$, and explain why two legs of equal vega but different implied vols do not have equal gamma.
\end{interviewq}
